\begin{frame}
    \setbeamertemplate{block begin}{
        \begin{beamercolorbox}[
                rounded=true,
                shadow=true,
            ]{block body}
        }
        \setbeamertemplate{block end}{
        \end{beamercolorbox}
    }
    \setbeamercolor{block body}{bg=zzu!5}

    \begin{block}{}
        \begin{minipage}[c]{0.13\textwidth}
            \centering
            \tikz\node[fill=zzu,minimum size=1.2cm,text=white,rounded corners=8pt]{\Large 01};
        \end{minipage}
        \begin{minipage}[c]{0.86\textwidth}
            \vspace{0.1cm}
            \textbf{\color{zzu}\large 构建 XX 框架 \& 设计 XX 模型}\\[0.3em]
            \small 构建 XX 框架并引入 XX 机制设计了一种融合 XX、XX、XX 与 XX 的四层神经网络架构并使用其在 XX 阶段做增强操作。
            \vspace{0.1cm}
        \end{minipage}
    \end{block}

    \vspace{0.2cm}

    \begin{block}{}
        \begin{minipage}[c]{0.13\textwidth}
            \centering
            \tikz\node[fill=zzu,minimum size=1.2cm,text=white,rounded corners=8pt]{\Large 02};
        \end{minipage}
        \begin{minipage}[c]{0.86\textwidth}
            \vspace{0.1cm}
            \textbf{\color{zzu}\large 开发算法实验平台并开源}\\[0.3em]
            \small 集成 XXX、XXX、XXX 大功能，本论文所用的全部代码已开源至：
            \newline
            \blueurl{https://github.com/XX/XX-Enhanced-XX}。
            \vspace{0.1cm}
        \end{minipage}
    \end{block}

    \vspace{0.2cm}

    \begin{block}{}
        \begin{minipage}[c]{0.13\textwidth}
            \centering
            \tikz\node[fill=zzu,minimum size=1.2cm,text=white,rounded corners=8pt]{\Large 03};
        \end{minipage}
        \begin{minipage}[c]{0.86\textwidth}
            \vspace{0.1cm}
            \textbf{\color{zzu}\large 创建多维验证体系}\\[0.3em]
            \small 结合对比实验和消融实验，通过 XX 指标评估 XX（XX、XX、XX 等），结合 Token 消耗与计算效率分析资源利用率。
            \vspace{0.1cm}
        \end{minipage}
    \end{block}
\end{frame}
